\subsubsection{Validity Checks and Robustness}

The reliability of the results obtained from the regression is based on the validity of the underlying assumptions of the regression model.
As the data is obtained from an observational study, where treatment is not randomly assigned, but based on the allocation of land through settlement schemes, potential confounding variables must be controlled.
Several validity and robustness checks are performed and will be described in the following sections.

First, a Covariate Balance Test is conducted to check the local randomization assumption of the RDD.
If allocation of settlement schemes is random around the cutoff, then the pre-existing conditions and baseline properties at the border should transition without any discontinuity from one side to the other.
To verify this assumption, the test is performed for two sets of properties, topographic and historical anthropogenic features.
For topographic features, time-invariant features of the land, specifically elevation and slope obtained from the SRTM dataset, are used as dependent variables.
Human settlement should not change these features, resulting in the $\tau$ of the regression being insignificant.
For historical anthropogenic features, the shortest distance from each grid cell to a historic city center and the distance to the nearest colonial railroad are used as dependent variables.
Testing for historical infrastructure features rules out the possibility of a selection bias driven by the proximity to existing markets.
If settlement schemes had been placed near those infrastructure features, differences in vegetation health could be driven by prolonged human activity rather than the settlement schemes themselves.
A significant $\tau$ would indicate this systematic difference in pre-existing conditions between the treatment and control group.
In contrast to other statistical tests, including covariates into the RDD cannot be used to fix a broken study design.
\textcite{cattaneoCovariateAdjustmentRegression2023} show that adding unbalanced covariates to the regression does not robustly restore the identification of the RDD,
as restrictive and not testable assumptions about the data generation process would be required.
As a result, the covariates used to improve the estimation strategy as described in the previous section are only used to increase the precision of the estimator.

Second, from an ecological perspective, the MSE-optimal bandwidth could be too narrow to capture the full effect of settlement schemes on vegetation indices.
Factors such as microclimate and change through human settlement are likely to influence the vegetation indices in neighboring grid cells.
To ensure that the estimated treatment effect is not only valid in the immediate vicinity of the cutoff line, a robustness check is performed using two alternative bandwidths.
For this purpose, larger bandwidths of 3000 and 9000 meters are used to check the estimator for a range of 10 and 30 times the size of one grid cell.

Furthermore, treating settlement scheme borders as ground truth for treatment entails the risk of violating the Stable Unit Treatment Value Assumption (SUTVA),
which requires that the treatment status of one unit does not affect the outcome of another unit.
Such spatial spillover effects could occur for two reasons:
First, settlers might have changed the vegetation of neighboring land outside the settlement scheme, for example by cutting trees for firewood or grazing livestock (encroachment),
which will be covered in more detail in the following chapter.
While completely enclosed settlement schemes were already excluded from the control group, grid cells located close to the cutoff line could still be affected by the treatment.
This would lead to an underestimation of the true effect of settlement schemes, as the vegetation indices of neighboring grid cells outside the settlement scheme would be reduced.
Second, historical maps of settlement scheme placement and their georeferencing could be inaccurate, and the plots actually allocated to settlers could differ from the official borders.
To solve both issues, a \textit{Donut RDD} robustness check is performed by excluding all grid cells located within 300 meters of the cutoff line.

Finally, assuming that the approximated borders of the settlement schemes and therefore the cutoff lines of non-georeferenced settlement schemes are accurate, could introduce spatial measurement errors.
As detailed in \cref{sec:data}, the area of about 30\% of settlement schemes could not be reconstructed from historical maps and are therefore rendered as circles,
calculated based on the allocated area and a point location of the center of the scheme.
This inaccuracy could lead to a misclassification of treatment status for grid cells located close to the cutoff line, which will influence the estimated treatment effect.
To rule this possibility out, a final robustness check is performed by excluding all grid cells from treatment and control groups associated with these settlement schemes from the regression.